% !TeX root = ../report.tex
Deep neural networks predict how cancer cells respond to drugs with high accuracy, but they are black boxes, and clinicians cannot easily tell whether a prediction rests on real biology or on artefacts in the data. Short IF--THEN rules extracted from a network can explain it, but only if the network itself is sound. This study tested whether training a drug-response network with Sharpness-Aware Minimisation (SAM), which steers training toward flatter minima and a smaller generalisation gap, produces better rules. We trained the DeepCCDS network on Genomics of Drug Sensitivity in Cancer data in three ways (Adam baseline, Adam with weight decay, and SAM with weight decay) and developed Per Drug Evolutionary Rule Extraction (PDERE), in a genetic-algorithm and an NSGA-II variant, to extract rules separately for 379 drugs. Rules were scored on held-out cell lines for fidelity, length and stability, ranked with PROMETHEE II, and compared across networks with paired Wilcoxon tests with false-discovery-rate control. SAM with weight decay cut the generalisation gap by 86\% (0.337 to 0.048), but mainly by fitting the training data less tightly, and validation error rose slightly. Rule quality did not follow the gap. Fidelity differences were at most about 0.014 on a 0 to 1 scale and were explained by weight decay, SAM alone made no detectable difference to fidelity, stability or PDERE-GA rule length, and known drug--gene links were almost never recovered. With one network per condition, this is a null result under a limited design, not evidence that SAM cannot help. We conclude that training metrics such as the generalisation gap are poor proxies for explanation quality, and that explanations must be audited directly.